\documentclass[11pt]{article}
\usepackage[margin=1in]{geometry}
\usepackage{amsmath, amssymb, amsthm}
\usepackage{tikz}
\usepackage{xcolor}
\usepackage{enumitem}
\usepackage{fancyhdr}

% Define brand colors
\definecolor{Garnet}{RGB}{115, 0, 10}
\definecolor{Atlantic}{RGB}{70, 106, 159}
\definecolor{Congaree}{RGB}{31, 65, 77}
\definecolor{NinetyBlack}{RGB}{54, 54, 54}
\definecolor{FiftyBlack}{RGB}{162, 162, 162}

\usetikzlibrary{arrows.meta, positioning, shapes.geometric}

\pagestyle{fancy}
\fancyhf{}
\rhead{J.C. Vaught}
\lhead{CSCE 775: Reinforcement Learning}
\cfoot{\thepage}

\newtheorem{problem}{Problem}
\theoremstyle{definition}
\newtheorem*{solution}{Solution}

\title{\textbf{CSCE 775: Reinforcement Learning} \\ Practice Quiz \\ Chapters 1, 3, 4}
\author{J.C. Vaught}
\date{Spring 2026}

\begin{document}

\maketitle

\section*{Instructions}
This quiz covers material from Chapters 1, 3, and 4. Show all work and justify your answers.

\newpage

\begin{problem}[Computing Returns with Discount Factors]
Given a discount factor $\gamma = 0.9$ and a reward sequence: $R_1 = 2$, $R_2 = -1$, $R_3 = 3$, $R_4 = 5$, $R_5 = 0$ with terminal time $T = 5$:
\begin{enumerate}[label=(\alph*)]
    \item Compute $G_0, G_1, G_2, G_3, G_4, G_5$ by working backwards from $G_5 = 0$.
    \item If $\gamma = 0.5$, what are $G_0$ through $G_5$?
\end{enumerate}
\end{problem}

\vfill
\newpage

\begin{problem}[Discounted Return for Infinite Sequence]
Given $\gamma = 0.8$ and a constant reward of $+3$ at every time step (continuing task):
\begin{enumerate}[label=(\alph*)]
    \item What is the return $G_t$ for any time $t$?
    \item If the first reward is $+10$ and all subsequent rewards are $+3$, what is $G_0$?
\end{enumerate}
\end{problem}

\vfill
\newpage

\begin{problem}[Bellman Equation Verification on Small MDP]
Consider a 3-state MDP with states $\{A, B, C\}$ and actions $\{a_1, a_2\}$:

\textbf{Transition dynamics:}
\begin{itemize}
    \item From $A$:
    \begin{itemize}
        \item $a_1$: go to $B$ with prob 0.7 (reward $+2$) or stay at $A$ with prob 0.3 (reward $+1$)
        \item $a_2$: go to $C$ with prob 1.0 (reward $+5$)
    \end{itemize}
    \item From $B$:
    \begin{itemize}
        \item $a_1$: go to $A$ with prob 0.4 (reward $+3$) or $C$ with prob 0.6 (reward $-1$)
        \item $a_2$: stay at $B$ with prob 1.0 (reward $+2$)
    \end{itemize}
    \item From $C$:
    \begin{itemize}
        \item $a_1$: go to $A$ with prob 1.0 (reward $+4$)
        \item $a_2$: go to $B$ with prob 0.5 (reward $+1$) or $C$ with prob 0.5 (reward $-2$)
    \end{itemize}
\end{itemize}

\textbf{Policy:} $\pi(a_1|A) = 0.7$, $\pi(a_2|A) = 0.3$; $\pi(a_1|B) = 0.5$, $\pi(a_2|B) = 0.5$; $\pi(a_1|C) = 0.6$, $\pi(a_2|C) = 0.4$

With $\gamma = 0.9$:
\begin{enumerate}[label=(\alph*)]
    \item Draw the transition diagram.
    \item Set up the Bellman equations for $v_\pi$.
    \item Solve the system of linear equations for $v_\pi(A)$, $v_\pi(B)$, $v_\pi(C)$.
\end{enumerate}
\end{problem}

\vfill
\newpage

\begin{problem}[Policy Evaluation on a Gridworld]
Consider a $3 \times 3$ gridworld with the terminal state at the top-right corner. Reward $= -1$ on all transitions. $\gamma = 1.0$ (undiscounted episodic task).

\begin{enumerate}[label=(\alph*)]
    \item Draw the gridworld with TikZ.
    \item Using a uniform random policy (all 4 actions equally likely), perform 3 iterations of iterative policy evaluation by hand.
    \item Start with $V(s) = 0$ for all states.
    \item Actions that would take the agent off the grid keep the agent in place.
    \item Show the value table after each sweep.
\end{enumerate}
\end{problem}

\vfill
\newpage

\begin{problem}[Policy Improvement]
Using the converged values from Problem 4, determine the greedy policy for each state. For each non-terminal state, show which action(s) maximize $q_\pi(s,a)$. Draw the resulting greedy policy with arrows.
\end{problem}

\vfill
\newpage

\begin{problem}[Value Iteration Steps]
Using the same $3 \times 3$ gridworld from Problem 4, apply value iteration:
\[
V(s) \leftarrow \max_a \sum_{s',r} p(s',r|s,a)[r + \gamma V(s')]
\]

Perform 3 sweeps of value iteration starting from $V(s) = 0$ for all states. Show the value table after each sweep and determine the optimal policy from the final values.
\end{problem}

\vfill
\newpage

\begin{problem}[Identifying Optimal Policies]
Consider a 2-state MDP with states $\{s_1, s_2\}$ and $\gamma = 0.9$:

\textbf{Transition dynamics:}
\begin{itemize}
    \item From $s_1$: action `stay' $\to s_1$ with reward $+1$; action `go' $\to s_2$ with reward $+5$
    \item From $s_2$: action `stay' $\to s_2$ with reward $+2$; action `go' $\to s_1$ with reward $-1$
\end{itemize}

\begin{enumerate}[label=(\alph*)]
    \item Write the Bellman optimality equations for $v_*(s_1)$ and $v_*(s_2)$.
    \item Solve for $v_*(s_1)$ and $v_*(s_2)$.
    \item Determine $\pi_*(s_1)$ and $\pi_*(s_2)$.
\end{enumerate}
\end{problem}

\vfill
\newpage

\begin{problem}[Conceptual Questions]
Answer the following:
\begin{enumerate}[label=(\alph*)]
    \item Why is the Bellman equation for $v_*$ nonlinear while the Bellman equation for $v_\pi$ is linear?
    \item What does it mean for a state to have the Markov property?
    \item Explain the difference between the reward signal and the value function. Why do we need both?
    \item In GPI, the evaluation and improvement processes ``compete.'' Explain what this means.
    \item Why might we prefer $q_*$ over $v_*$ for determining optimal actions?
\end{enumerate}
\end{problem}

\vfill
\newpage

\begin{problem}[Computing $q_\pi$ from $v_\pi$]
Using the 3-state MDP from Problem 3 and the computed $v_\pi$ values:
\begin{enumerate}[label=(\alph*)]
    \item Compute $q_\pi(A, a_1)$, $q_\pi(A, a_2)$, $q_\pi(B, a_1)$, $q_\pi(B, a_2)$
    \item Verify that $v_\pi(A) = \sum_a \pi(a|A) q_\pi(A, a)$
\end{enumerate}

Use the values: $v_\pi(A) \approx 31.85$, $v_\pi(B) \approx 25.74$, $v_\pi(C) \approx 30.21$ (from Problem 3 solution).
\end{problem}

\vfill
\newpage

\begin{problem}[Action-Value Bellman Optimality]
For the 2-state MDP from Problem 7:
\begin{enumerate}[label=(\alph*)]
    \item Write the Bellman optimality equations for $q_*(s,a)$ for all state-action pairs.
    \item Solve for all $q_*$ values.
    \item Verify that $v_*(s) = \max_a q_*(s,a)$ for both states.
\end{enumerate}
\end{problem}

\vfill
\newpage

\end{document}
